\problema{Vertical Order Traversal of a Binary Tree}{987}{src/02-arboles/987-vertical-order.cpp}{H}

A cada nodo del árbol le asignamos una \emph{columna} y una \emph{fila}: la
raíz vive en columna \texttt{0}, fila \texttt{0}; el hijo izquierdo se mueve
una columna a la izquierda y una fila hacia abajo; el hijo derecho, una
columna a la derecha y una fila hacia abajo. Debemos devolver los valores
agrupados de columna menor a mayor; dentro de cada columna, ordenados por
fila ascendente; y si dos nodos caen exactamente en la misma columna y la
misma fila, se desempatan por \textbf{valor ascendente}. Este último detalle
es justo lo que distingue a este problema del más sencillo \#314 y lo que lo
hace Hard: no basta con recorrer el árbol, hay que ordenar con un criterio
compuesto de tres niveles.

\paragraph{Intuición.} Si en vez de pensar en "izquierda/derecha" pensamos en
coordenadas, cada nodo termina siendo una terna
$(\text{columna}, \text{fila}, \text{valor})$. Una vez que tenemos todas las
ternas, el problema se reduce a \emph{ordenarlas}: por columna, luego por
fila, luego por valor. El orden natural de una tupla ya hace exactamente eso.
No hace falta ninguna estructura de datos sofisticada: un DFS que acumula
ternas y un \texttt{sort} al final resuelven el problema completo.

\begin{center}
\begin{verbatim}
          3
        /   \
       9     20
            /  \
          15    7
\end{verbatim}
\end{center}

\begin{center}
\begin{tabular}{lll}
\textbf{columna} & \textbf{fila} & \textbf{valor} \\
$-1$ & 1 & 9 \\
$0$  & 0 & 3 \\
$0$  & 2 & 15 \\
$1$  & 1 & 20 \\
$2$  & 2 & 7 \\
\end{tabular}
\end{center}

Al agrupar por columna (ya ordenado) obtenemos
\texttt{[[9],[3,15],[20],[7]]}: exactamente el resultado esperado.

\paragraph{Solución.} Un DFS (o BFS, es indistinto) recorre el árbol y por
cada nodo guarda la terna \texttt{(columna, fila, valor)} en un vector. El
hijo izquierdo se visita con \texttt{columna-1, fila+1}; el derecho con
\texttt{columna+1, fila+1}. Al terminar, ordenamos el vector de ternas: como
\texttt{tuple<int,int,int>} ya compara componente a componente, un solo
\texttt{sort} deja todo en el orden pedido, incluido el desempate por valor.
Por último recorremos el vector ordenado agrupando los valores consecutivos
que comparten columna.

El caso donde el desempate por valor realmente importa: en el árbol completo
\texttt{[1,2,3,4,5,6,7]}, el nodo \texttt{5} (hijo derecho de \texttt{2}) y el
nodo \texttt{6} (hijo izquierdo de \texttt{3}) caen ambos en columna \texttt{0},
fila \texttt{2}. Sin el desempate por valor el orden entre ellos sería
ambiguo; con él, \texttt{5} siempre va antes que \texttt{6}.

\lstinputlisting{src/02-arboles/987-vertical-order.cpp}

\complejidad{$O(n \log n)$}{$O(n)$}

\paragraph{Notas de entrevista (Amazon).} El error típico es asumir que un
recorrido en orden (BFS por niveles o DFS in-order) ya entrega el resultado
final: no es así, porque nodos de columnas distintas pueden intercalarse en
el mismo nivel del recorrido. La forma robusta es separar completamente las
dos fases: primero \emph{recolectar} toda la información posicional de cada
nodo (columna, fila, valor) sin preocuparse por el orden, y después
\emph{ordenar} con el criterio compuesto correcto. Vale la pena mencionar en
la entrevista que este problema es la versión "difícil" de LeetCode 314
(Binary Tree Vertical Order Traversal): la diferencia entera está en el
desempate por valor cuando coinciden columna y fila, algo fácil de pasar por
alto si se resuelve de memoria.
